\section{Related Work}\label{sec:related}

\scl{} is related to four lines of research: the trackers it works with, the
confidence of the detectors that feed them, trackers that adapt their own
thresholds, and methods that adapt inference while running. They are
reviewed below in this order, and Table~\ref{tab:related} compares the
closest methods with \scl{}.

Most practical trackers follow the tracking-by-detection design. SORT
predicts each box with a Kalman filter and matches the predictions to new
detections by overlap with the Hungarian method \cite{bewley2016sort}.
ByteTrack added a second matching step that uses low-score detections to
continue existing tracks \cite{zhang2022bytetrack}, and many later trackers
kept this two-stage design. Others, such as OC-SORT \cite{cao2023ocsort}, use
only one score threshold. All these trackers read detector scores through
fixed thresholds that were tuned for one detector, usually YOLOX
\cite{ge2021yolox}. Seven published trackers of both kinds are used in this
paper as hosts and are not changed (Sect.~\ref{sec:setup}).

These fixed thresholds are a weak point, because scores do not mean the same
thing for every detector. Detectors of different families, such as YOLOX
\cite{ge2021yolox}, YOLOv8 \cite{jocher2023yolov8}, YOLOv9
\cite{wang2024yolov9}, DETR \cite{carion2020detr}, and RT-DETR
\cite{zhao2024rtdetr}, give different score distributions. Their confidence
is not calibrated across models and conditions, and calibration methods for
object detection learn a mapping from labelled data
\cite{kuppers2020detcalib}. A tracker tuned for one detector therefore cannot
assume that another detector's scores mean the same thing, and labels to
recalibrate them are rarely available after deployment. Image processing
offers a tool that needs no labels: Otsu's method picks a threshold from a
histogram so that the two resulting groups are as separate as possible
\cite{otsu1979}. \scl{} applies it to the logits (log-odds) of the detector
scores of recent frames.

Some trackers already adapt their own thresholds. Adaptive ByteTrack sets the
score threshold of ByteTrack from the detections of each frame
\cite{ma2023adaptivebytetrack}. AdapTrack adapts the thresholds of its
matching step \cite{shim2024adaptrack}. ACT sets the high and low thresholds
of each frame from the first quartile of the scores and the number of
candidates \cite{act2026adaptive}. These methods are built into one tracker
and use only the current frame. \scl{} differs in four ways. It uses past
frames. It decides for the whole stream whether to act at all, and otherwise
keeps the tracker's own settings. It works through a declared host contract
on several published trackers without changing them. And it is tested as one
frozen configuration under score changes and on external systems.

A wider group of methods adapts inference itself. Dynamic neural networks
adapt their computation to each input \cite{han2022dynamic}, and Chameleon
re-profiles the settings of video analytics pipelines over time
\cite{jiang2018chameleon}. These methods learn or measure a policy for each
model. Our earlier work, AC-MOT \cite{ismail2026acmot}, is a controller that
needs no training and sets the detector operating point (score threshold,
NMS threshold, and input resolution) while running, from image and track
cues. Such a controller is tied to the detector it was calibrated with,
because a change in the detector's score scale does not change the image.
\scl{} reads the scores instead and acts on the tracker's input and
thresholds.

In summary, current methods either adapt thresholds inside one tracker from
the current frame, calibrate scores with labelled data, or learn a policy for
one model. What is missing is a layer that needs no training, adjusts itself
from the detector's own scores, works with many unchanged trackers, and keeps
a well-matched tracker at its own accuracy. \scl{} is designed to fill this
gap. Table~\ref{tab:related} summarizes the differences.

\begin{table}[t]
\caption{Summary of the adaptive methods closest to \scl{} (BT: ByteTrack).}
\label{tab:related}
\centering\footnotesize
\setlength{\tabcolsep}{3pt}
\begin{tabular}{@{}>{\raggedright\arraybackslash}p{17mm}>{\raggedright\arraybackslash}p{18mm}>{\raggedright\arraybackslash}p{15mm}>{\raggedright\arraybackslash}p{18mm}@{}}
\toprule
Method & Adapted & Statistics & Where\\
\midrule
Adapt.\ BT \cite{ma2023adaptivebytetrack} & confidence threshold & current frame & inside BT\\
AdapTrack \cite{shim2024adaptrack} & matching thresholds & current frame & inside the tracker\\
ACT \cite{act2026adaptive} & high and low thresholds & current frame & inside the tracker\\
K{\"u}ppers \cite{kuppers2020detcalib} & detector scores & labelled data & detector output\\
Chameleon \cite{jiang2018chameleon} & pipeline configuration & profiling & analytics pipeline\\
AC-MOT \cite{ismail2026acmot} & detector operating point & image and track cues & before the detector\\
\scl{} (this paper) & candidates, scores, thresholds & past frames, no labels & between detector and tracker\\
\bottomrule
\end{tabular}
\end{table}
